\documentclass{article}

\usepackage{amsmath,amssymb}
\usepackage{xurl}
\usepackage[hidelinks]{hyperref}
\setlength{\emergencystretch}{2em}

\input{command}

\title{The Nonvanishing Partner in Wolf's Kramers Theorem~II}
\author{}
\date{2026-08-24}

\begin{document}
\maketitle
\gapnote{false-source}{resolved}

\section{The assertion}

Wolf's Kramers theorem~II states:\footnote{\cite[Chapter~3, Theorem~3.2
(``Kramer's theorem~II'')]{Wolf2012Quantum}; local source
\path{Notes/WolfNoteTexSource/ch03_positive_not_completely.tex},
lines~527--530.  Tracked in \ghissue{17}.}
\begin{align}
  \left.
  \begin{minipage}{0.82\linewidth}
  Let $H$ be Hermitian and $HA=AH^{T}$ for some anti-symmetric $A\neq 0$.
  Then each eigenvalue of $H$ is at least two-fold degenerate.
  \end{minipage}
  \right\}
  \tag{\path{Wolf ch03, lines 527--530}}
\end{align}
The theorem is offered as a relaxation of the first Kramers theorem
(\cite[Chapter~3, Theorem~3.1]{Wolf2012Quantum}), which
treats a Hermitian $H$ commuting with an anti-unitary $T$ of square $-\Id$:
writing $T=\Gamma V$ with $\Gamma$ complex conjugation and $V$ unitary, the
hypotheses become $HV^{\dagger}=V^{\dagger}H^{T}$ and $V^{T}=-V$, and the
printed text adds that ``the restriction on the required symmetry in the
theorem can be relaxed as the same proof shows.''\footnote{Local source,
lines~498--501 (the matrix reduction) and lines~523--525 (the relaxation
claim).}

\section{The point at issue}

The proof in question produces, for every eigenvector
$H\ket{\psi}=\lambda\ket{\psi}$, the partner vector $A\ket{\overline\psi}$ and
argues that the two vectors are orthogonal eigenvectors for the same
eigenvalue.  Both calculation steps survive the relaxation:
\begin{align}
  H\,A\ket{\overline\psi}
  =A H^{T}\ket{\overline\psi}
  =\lambda\,A\ket{\overline\psi},
  \qquad
  \braket{\psi}{A\overline\psi}=0,
\end{align}
the first because Hermiticity makes $\lambda$ real and
$H^{T}\ket{\overline\psi}=\overline{H\ket{\psi}}$, the second because the
quadratic form of an anti-symmetric matrix is its own negative.  What does
\emph{not} survive is the nondegeneracy step.  In the first theorem the
partner is $V^{\dagger}\ket{\overline\psi}$ with $V$ unitary, hence nonzero;
for a general anti-symmetric $A\neq 0$ the vector $\ket{\overline\psi}$ may lie
in the kernel of $A$, and then the partner vanishes and the orthogonality
statement is vacuous.

Concretely, on $\C^{3}$ take
\begin{align}
  H=\operatorname{diag}(0,1,1),
  \qquad
  A=
  \begin{pmatrix}
    0&0&0\\
    0&0&1\\
    0&-1&0
  \end{pmatrix}.
\end{align}
Then $H=H^{\dagger}$, $A^{T}=-A$, $A\neq 0$, and $HA=AH^{T}=A$, since the
nontrivial rows and columns of $A$ sit entirely inside the $1$-eigenspace of
$H$.  The eigenvalue $0$ of $H$ is simple, with eigenvector $e_{1}$, and the
partner vanishes: $A\overline{e_{1}}=Ae_{1}=0$.  So the printed conclusion
fails while every printed hypothesis holds.

\section{The corrected statement}

The calculation above proves the conditional statement: the eigenvalue
$\lambda$ of an eigenvector $\ket{\psi}$ is at least two-fold degenerate
\emph{whenever the partner is nonzero},
\begin{align}
  A\ket{\overline\psi}\neq 0
  \Longrightarrow
  \dim\ker(H-\lambda)\ge 2,
\end{align}
because $\ket{\psi}$ and $A\ket{\overline\psi}$ are then two orthogonal,
hence linearly independent, eigenvectors for $\lambda$.  A hypothesis-free
global conclusion therefore requires $A$ to be injective, so that
$\ket{\psi}\neq 0$ implies $A\ket{\overline\psi}\neq 0$: if the
anti-symmetric intertwiner $A$ is invertible, then every eigenvalue of $H$ is
at least two-fold degenerate.  Over $\C$ an anti-symmetric matrix can be
invertible only in even dimensions, since
$\det A=\det A^{T}=(-1)^{d}\det A$; this matches the observation printed next
to the first theorem that anti-symmetric unitaries exist only in even
dimensions.\footnote{Local source, lines~523--524.}

The first Kramers theorem is exactly the invertible case: for $T=\Gamma V$
the choice $A=V^{\dagger}$ is anti-symmetric (because $V^{T}=-V$) and
invertible (because $V$ is unitary), so the corrected theorem~II specializes
to theorem~I, with Wolf's own proof.

\section{The formalization}

The corrected development lives in
\doclink{QICLean/Algebra/KramersDegeneracy}.\footnote{Transport:
\leanid{Matrix.IsHermitian.mulVec_star_intertwiner_of_mul_eq_mul_transpose};
orthogonality:
\leanid{Matrix.dotProduct_star_mulVec_star_eq_zero_of_transpose_eq_neg};
conditional theorem:
\leanid{Matrix.IsHermitian.two_le_finrank_eigenspace_of_intertwiner_mulVec_star_ne_zero};
invertible global corollary:
\leanid{Matrix.IsHermitian.two_le_finrank_eigenspace_of_antisymmetric_isUnit};
theorem~I in matrix form, derived as the case $A=V^{\dagger}$:
\leanid{Matrix.IsHermitian.two_le_finrank_eigenspace_of_antisymmetric_unitary};
theorem~I with the printed anti-unitary hypotheses:
\leanid{Matrix.IsHermitian.two_le_finrank_eigenspace_of_antiunitary}.}
Degeneracy is recorded as $2\le\dim$ of the eigenspace; for a Hermitian
matrix the geometric and algebraic multiplicities agree, so nothing is lost
against the printed phrase ``two-fold degenerate''.  The false printed
statement itself is deliberately not formalized.

\section{Conclusion}

Kramers theorem~II as printed is false; the hypothesis $A\neq0$ is
insufficient because the partner $A\ket{\overline\psi}$ can vanish.  The
correct local conclusion holds under the exact nonvanishing condition
$A\ket{\overline\psi}\neq0$, and the global conclusion holds when $A$ is
invertible.  This is a genuine source error, resolved by the corrected
statements above.

\bibliographystyle{alpha}
\bibliography{references}

\end{document}
